%========================================================================
% DOCUMENTO COMPILABLE AUTÓNOMO
% Apuntes de Inglés Jurídico: El Poder Judicial de los Estados Unidos
% Método Cornell integrado
%========================================================================

\documentclass[12pt,oneside]{book}

%---- PAQUETES OBLIGATORIOS ----
\usepackage[utf8]{inputenc}
\usepackage[T1]{fontenc}
\usepackage[spanish,es-nodecimaldot]{babel}

%---- MÁRGENES ----
\usepackage[
    top=2.5cm,
    bottom=2.5cm,
    left=3cm,
    right=2.5cm,
    headheight=15pt
]{geometry}

%---- MATEMÁTICA ----
\usepackage{amsmath}
\usepackage{amssymb}
\usepackage{mathtools}

%---- IMÁGENES ----
\usepackage{graphicx}
\usepackage{float}
\usepackage{caption}
\usepackage{subcaption}

%---- TABLAS ----
\usepackage{booktabs}
\usepackage{array}
\usepackage{longtable}
\usepackage{tabularx}

%---- LISTAS ----
\usepackage{enumitem}

%---- COLORES Y CAJAS ----
\usepackage{xcolor}
\usepackage[most]{tcolorbox}

%---- ENCABEZADOS ----
\usepackage{fancyhdr}

%---- HIPERVÍNCULOS Y METADATOS PDF ----
\usepackage[
    colorlinks=true,
    linkcolor=blue!50!black,
    citecolor=green!40!black,
    urlcolor=blue!60!black,
    pdftitle={Apuntes de Inglés Jurídico: El Poder Judicial de los Estados Unidos},
    pdfauthor={Isaac Edgar Camacho Ocampo},
    pdfsubject={Apuntes universitarios},
    bookmarks=true
]{hyperref}

%---- CONFIGURACIÓN GLOBAL ----
\setcounter{tocdepth}{2}
\setlength{\parindent}{0pt}
\setlength{\parskip}{0.5em}

\setlist[itemize]{
    topsep=0.3em,
    itemsep=0.2em
}

\setlist[enumerate]{
    topsep=0.3em,
    itemsep=0.2em
}

\clubpenalty=10000
\widowpenalty=10000

%---- RUTAS DE IMÁGENES ----
\graphicspath{
    {./img/}
    {./graficos/}
    {./figuras/}
}

%---- ENCABEZADOS Y PIES ----
\pagestyle{fancy}
\fancyhf{}

\fancyhead[L]{\nouppercase{\leftmark}}
\fancyhead[R]{Inglés Jurídico}

\fancyfoot[C]{\thepage}

\renewcommand{\headrulewidth}{0.4pt}
\renewcommand{\footrulewidth}{0pt}

\fancypagestyle{plain}{
    \fancyhf{}
    \fancyfoot[C]{\thepage}
    \renewcommand{\headrulewidth}{0pt}
}

%---- COMANDOS MATEMÁTICOS ----
\newcommand{\abs}[1]{\left|#1\right|}
\newcommand{\norm}[1]{\left\lVert#1\right\rVert}
\newcommand{\vect}[1]{\mathbf{#1}}
\newcommand{\deriv}[2]{\frac{d#1}{d#2}}
\newcommand{\pderiv}[2]{\frac{\partial#1}{\partial#2}}

%---- CAJAS ACADÉMICAS ----
\newtcolorbox{definition}[1]{
    enhanced,
    breakable,
    title={Definición: #1},
    fonttitle=\bfseries,
    colback=blue!3,
    colframe=blue!50!black,
    boxrule=0.8pt,
    arc=2mm
}

\newtcolorbox{important}{
    enhanced,
    breakable,
    title={Importante},
    fonttitle=\bfseries,
    colback=yellow!8,
    colframe=orange!70!black,
    boxrule=0.8pt,
    arc=2mm
}

\newtcolorbox{example}[1]{
    enhanced,
    breakable,
    title={Ejemplo: #1},
    fonttitle=\bfseries,
    colback=green!4,
    colframe=green!40!black,
    boxrule=0.8pt,
    arc=2mm
}

\newtcolorbox{formula}[1]{
    enhanced,
    breakable,
    title={Fórmula / Regla: #1},
    fonttitle=\bfseries,
    colback=purple!4,
    colframe=purple!50!black,
    boxrule=0.8pt,
    arc=2mm
}

\newtcolorbox{exercise}[1]{
    enhanced,
    breakable,
    title={Ejercicio: #1},
    fonttitle=\bfseries,
    colback=gray!5,
    colframe=gray!60!black,
    boxrule=0.8pt,
    arc=2mm
}

\newtcolorbox{note}{
    enhanced,
    breakable,
    title={Nota},
    fonttitle=\bfseries,
    colback=white,
    colframe=gray!50,
    boxrule=0.6pt,
    arc=2mm
}

%---- METADATOS DEL DOCUMENTO ----
\newcommand{\universidad}{Universidad de Buenos Aires}
\newcommand{\carrera}{Carrera de Traductorado en Ámbito Jurídico}
\newcommand{\materia}{Inglés Jurídico --- Lectocomprensión}
\newcommand{\titulo}{El Poder Judicial de los Estados Unidos}
\newcommand{\autor}{Isaac Edgar Camacho Ocampo}
\newcommand{\anio}{2026}

%========================================================================
\begin{document}

%====== PORTADA ======
\begin{titlepage}

\thispagestyle{empty}

\begin{center}

\vspace*{1cm}

{\large \textsc{\universidad}}

\vspace{0.2cm}

{\large \carrera}

\vspace{3cm}

{\Huge \textbf{\materia}}

\vspace{1cm}

{\LARGE \titulo}

\vspace{0.8cm}

\textit{Comprender la estructura judicial es fundamental para interpretar correctamente textos y conflictos en el ámbito del derecho internacional.}

\vspace{1.2cm}

\rule{0.70\textwidth}{0.5pt}

\vspace{0.5cm}

{\large Apuntes de estudio}

\vspace{0.5cm}

\rule{0.70\textwidth}{0.5pt}

\vfill

{\large \autor}

\vspace{0.3cm}

{\large \anio}

\end{center}

\vfill

\begin{flushleft}
\textit{Este es un modesto aporte a los estudiantes de ``Inglés Jurídico'' no pretende sustituir el material oficial de la materia.}
\end{flushleft}

\end{titlepage}

%====== ÍNDICE ======
\tableofcontents

\newpage

%====== INTRODUCCIÓN ======
\chapter*{Introducción}

\section*{¿De qué trata esta clase?}

La presente clase trabaja sobre un texto oficial en inglés que describe la organización del Poder Judicial federal de los Estados Unidos: la Corte Suprema, los tribunales federales inferiores y las agencias judiciales especializadas. A partir de ese texto se entrena una estrategia de lectocomprensión mediante la identificación de términos transparentes, el análisis de la macroestructura y la activación de saberes previos.

En paralelo, se profundiza en cuestiones lingüísticas propias del inglés jurídico: polisemia, falsos cognados (\textit{court}, \textit{sentence}, \textit{jurisdiction}), verbos modales (\textit{must}, \textit{may}, \textit{might}, \textit{should}, \textit{shall}, \textit{would}, \textit{could}) y phrasal verbs. Constantemente se compara el sistema estadounidense con el sistema judicial argentino.

\section*{¿Para qué sirve esto en la vida profesional?}

\begin{itemize}
    \item \textbf{Para el ejercicio profesional del derecho y la traducción:} quien trabaje con contratos, sentencias, jurisprudencia comparada o clientes extranjeros necesita leer con precisión textos jurídicos en inglés. Confundir \textit{sentence} con \textit{sentencia} o \textit{jurisdiction} con \textit{jurisdicción} no es un matiz estético: puede cambiar por completo el sentido de una cláusula, una pena o un documento legal.
    
    \item \textbf{Para la traducción pública y pericial:} identificar correctamente los verbos modales (\textit{must}, \textit{shall}, \textit{should}) permite distinguir entre una obligación legal, una recomendación y una facultad --- una distinción central en contratos y normas, donde de eso depende la validez y el alcance de lo que se traduce.
    
    \item \textbf{Para el derecho comparado:} entender cómo se organiza el Poder Judicial de otro país (instancias, competencias, jurado, procedimiento) entrena una habilidad transferible a cualquier sistema jurídico extranjero, cada vez más necesaria en un mundo de relaciones internacionales, arbitrajes y litigios transfronterizos.
    
    \item \textbf{Para la vida cotidiana y la formación general:} aprender a leer con estrategia --- apoyándose en el contexto, en la estructura del texto y en los saberes previos, en lugar de traducir palabra por palabra --- es una herramienta de estudio y de pensamiento crítico útil mucho más allá de esta materia.
\end{itemize}

\newpage

%====== PARTE 1: HERRAMIENTAS LINGÜÍSTICAS ======
\chapter{Herramientas Lingüísticas para el Inglés Jurídico}

\section{Polisemia: los múltiples sentidos de ``Law''}

El término \textbf{\textit{law}} (ley) es un término polisémico: tiene más de un significado. Por contexto nos damos cuenta a qué se refiere. Los significados no son independientes unos de otros, pero es fundamental identificar cuál corresponde según la situación.

\subsection*{Significados principales de ``law''}

\begin{table}[H]
\centering
\begin{tabularx}{\textwidth}{p{3.5cm} X p{4.5cm}}
\toprule
\textbf{Significado} & \textbf{Explicación} & \textbf{Ejemplo} \\
\midrule
Una norma en particular & Se refiere a una ley específica dentro de un sistema jurídico. & \textit{the law of business associations} (la ley de sociedades comerciales) \\
\midrule
Un sistema jurídico & El conjunto de normas, decisiones, fallos, decretos y principios jurídicos de un ordenamiento. & \textit{a determined legal system} \\
\midrule
Un área o rama del derecho & Una materia específica dentro de la profesión. & \textit{family law}, \textit{business law}, \textit{criminal law} (derecho de familia, comercial, penal) \\
\midrule
La profesión / las fuerzas de seguridad & De forma más informal, alude al conjunto de jueces, abogados y fiscales, o incluso a la policía. & ``ahí viene la ley'' (viene la policía) \\
\bottomrule
\end{tabularx}
\end{table}

Cuando el texto original señala que la corte ``applies the law to individual cases'', se está hablando de la primera acepción: aplica la norma general a un caso particular.

\section{Falsos Cognados I: Court, Sentence y Jurisdiction}

\subsection{``Court'': no siempre es ``corte''}

La traducción de \textit{court} es engañosa si se traduce siempre como ``corte''. 

\textit{Court} es el término genérico para tribunal. Sin embargo:

\begin{itemize}
    \item ``Corte'' solo corresponde a la Corte Suprema.
    \item Un tribunal de alzada es una \textbf{cámara}.
    \item Un tribunal de primera instancia es un \textbf{juzgado} o \textbf{tribunal}.
\end{itemize}

Traducir siempre ``corte'' es un error frecuente, sobre todo si se trata de un tribunal de primera instancia.

\begin{example}{Superior Court en Nueva York}
Un tribunal identificado como ``Superior Court'' puede resultar engañoso. Un divorcio de trámite común no llegaría, en el sistema comparado, a una instancia de alzada. Ese ``Superior Court'' era, en realidad, un tribunal de primera instancia. Ante nombres de ese tipo conviene consultar el organigrama judicial oficial de cada estado antes de traducir.
\end{example}

\subsection{``Sentence'': un falso cognado de importancia crítica}

El término \textit{sentence} es un falso cognado especialmente problemático.

\begin{definition}{Sentence vs. Sentencia}
En el lenguaje cotidiano, \textit{sentence} significa ``oración''. En el lenguaje jurídico, \textit{sentence} se refiere específicamente a la \textbf{pena} impuesta a quien ya fue declarado culpable (prisión, multa, pena de muerte, etc.). Solo existe cuando hubo condena.

La palabra española ``sentencia'' ---que puede ser condenatoria o absolutoria--- se traduce al inglés como \textbf{\textit{judgement}} o \textbf{\textit{resolution}}.
\end{definition}

La \textit{US Sentencing Commission} es la comisión federal que fija pautas y lineamientos uniformes para la imposición de penas, no para su ejecución ni para la determinación de la responsabilidad.

\subsection{``Jurisdiction'': jurisdicción y competencia en un solo término}

En Argentina, \textbf{jurisdicción} se refiere a la facultad o potestad que tienen todos los jueces por su investidura, por tener su cargo. La \textbf{competencia}, en cambio, es la medida de la jurisdicción: es cómo esa jurisdicción se distribuye entre distintas autoridades en razón de la materia, de la persona, del territorio, del monto, del grado.

\begin{important}
La conclusión práctica para la lectura de textos jurídicos en inglés: 
\begin{itemize}
    \item \textbf{\textit{Jurisdiction}} en textos de doctrina = ``jurisdicción''
    \item \textbf{\textit{Jurisdiction}} en la práctica (demandas, fallos, textos procesales) = ``competencia''
\end{itemize}
\end{important}

Cuando se dice que la Corte Suprema tiene ``appellate jurisdiction'' en la mayoría de las causas, se está hablando de su competencia por apelación. Cuando en algunos casos excepcionales tiene ``original jurisdiction'', se habla de su competencia originaria.

\section{Comparativos y Superlativos en el Lenguaje Jurídico}

El comparativo compara un elemento A con otro elemento (uno o dos). El superlativo compara ese elemento con todo el universo de elementos similares.

\begin{example}{Comparativo vs. Superlativo}
Si comparás a una persona con otra y decís ``es más alto que fulano'', estás comparando A con B (comparativo). Cuando decís ``the highest'', lo que haces es comparar A con B, con C, con D, con todo lo que existe de ese universo. El artículo ``the'' ya indica que estás comparando con todo el universo de cosas similares.
\end{example}

\subsection*{Aplicación al sistema judicial}

\begin{table}[H]
\centering
\begin{tabularx}{\textwidth}{p{3cm} p{2.5cm} X}
\toprule
\textbf{Término} & \textbf{Tipo} & \textbf{Uso} \\
\midrule
\textbf{highest} & Superlativo & La Corte Suprema es el tribunal más alto: no hay ningún tribunal por encima de ella. \\
\midrule
\textbf{higher / lower} & Comparativos & Se usan para comparar dos instancias entre sí (por ejemplo, el tribunal de alzada frente al tribunal inferior cuyo fallo se apela). \\
\midrule
\textbf{upper} & Comparativo & Compara un tribunal superior con otro que no lo es; no implica que sea el más alto del sistema. \\
\bottomrule
\end{tabularx}
\end{table}

Estos comparativos son siempre relativos al elemento con el que se los compara: un tribunal puede ser ``inferior'' a la Corte Suprema sin ser necesariamente de primera instancia (por ejemplo, una cámara de apelaciones es ``inferior'' a la Corte Suprema, pero superior a un tribunal de primera instancia).

\section{Verbos Modales y el Lenguaje del Deber Ser}

Los verbos modales son auxiliares: no quieren decir nada por sí solos. Necesitan ir acompañados de otro verbo, que es el que da sentido a la acción. El modo que indican afecta a ese otro verbo. En general, lo que estos verbos expresan es el deber ser, concepto central en el derecho.

\subsection*{Verbos modales principales}

\begin{table}[H]
\centering
\begin{tabularx}{\textwidth}{p{2.5cm} X p{4.5cm}}
\toprule
\textbf{Modal} & \textbf{Idea que expresa} & \textbf{Ejemplo} \\
\midrule
\textbf{can} & Habilidad / posibilidad concreta. & ``Can I go to the bathroom?'' --- remite a la capacidad de hacerlo. \\
\midrule
\textbf{may} & Permiso (más educado que \textit{can}) / posibilidad / otorgamiento de permiso. & ``May I go to the bathroom?'' (pido permiso); ``You may go'' (te doy permiso). \\
\midrule
\textbf{must} & Obligación. & ``There must be a quorum of six'' --- requisito ineludible para resolver una causa. \\
\midrule
\textbf{might} & Probabilidad remota o posibilidad más incierta que \textit{may}. & ``It might rain today.'' \\
\midrule
\textbf{should} & Recomendación o sugerencia (no obligación). & ``You should see a professor'' --- una sugerencia, no una obligación. \\
\midrule
\textbf{shall} & La obligación por excelencia; propia del lenguaje jurídico, contractual y religioso. & ``You shall not kill'' (mandamiento); ``The seller shall deliver the goods'' (contratos). \\
\midrule
\textbf{would / could} & Equivalentes al condicional español; pedidos educados. & ``Would you go?''; ``Could you please tell me?'' \\
\bottomrule
\end{tabularx}
\end{table}

\subsection{El caso de ``Shall'': el verbo del deber ser jurídico}

\textit{Shall} es la obligación por excelencia. Es raro verlo en un texto cotidiano que no tenga que ver con derecho o algo formal. En lo religioso aparece en los mandamientos (``you shall not kill'', ``no matarás''), donde hay una obligación moral que emana de una autoridad religiosa. En el caso del derecho, esas obligaciones marcadas con \textit{shall} están determinadas por una autoridad superior: la ley, la norma, los principios jurídicos.

\begin{example}{Shall en contratos}
``The purchaser shall pay the price'' (el comprador deberá abonar el precio). ``The seller shall deliver the goods'' (el vendedor deberá entregar los bienes). Aunque esa obligación surge de lo pactado entre las partes, en definitiva proviene de algo superior: la norma que reconoce que los contratos con objeto lícito son vinculantes.
\end{example}

En el lenguaje jurídico español, este uso de \textit{shall} se traduce frecuentemente como futuro con sentido de obligación: ``the parties shall agree'', ``las partes deberán acordar''. Ese uso del futuro con sentido de obligación es muy propio del lenguaje jurídico en español.

\newpage

%====== PARTE 2: LA CORTE SUPREMA ======
\chapter{La Corte Suprema de Justicia}

\section{Composición y Quórum de la Corte Suprema}

La Corte Suprema de los Estados Unidos está integrada por \textbf{nueve miembros}:

\begin{itemize}
    \item Un \textbf{Chief Justice} (Presidente de la Corte)
    \item Ocho \textbf{Associate Justices} (Jueces Asociados o Ministros)
\end{itemize}

\begin{definition}{Chief Justice}
El \textit{Chief Justice} es el presidente de la Corte Suprema. Aunque a menudo se traduce literalmente como ``justicia en jefe'', la expresión correcta es ``presidente de la Corte''. Se sienta en el centro y es la cabeza del Poder Judicial federal.
\end{definition}

Para resolver una causa (\textit{to decide a case}) es necesario un \textbf{quórum}. El texto señala: ``there must be a quorum of six justices''. Es decir, se necesitan \textbf{seis jueces presentes, mínimo}, de los nueve totales. El uso de \textit{must} indica que se trata de una obligación: si no están presentes los seis jueces, la Corte no puede resolver la causa.

\section{Mandato Vitalicio: Term y Serve}

El texto señala que los jueces no tienen un mandato con plazo fijo (``there is no fix term for justices''). La traducción literal resultaría forzada en español. Lo que se expresa es que \textbf{no tienen un mandato por un período determinado}: es de por vida, salvo que se jubilen, se mueran o los remuevan por algún motivo particular, con los motivos establecidos por ley.

\begin{table}[H]
\centering
\begin{tabularx}{\textwidth}{p{3cm} X p{3.5cm}}
\toprule
\textbf{Término} & \textbf{En este contexto} & \textbf{Otro significado} \\
\midrule
\textbf{term} & Mandato (duración del cargo de un juez). & Plazo de un contrato: ``the term of the agreement''. \\
\midrule
\textbf{serve} & Ejercer, cumplir, desempeñar una función (no ``servir'' en sentido literal). & Puede acercarse a ``prestar servicio'', según el contexto. \\
\bottomrule
\end{tabularx}
\end{table}

En el sistema estadounidense, mientras no medie jubilación, renuncia, fallecimiento o remoción, el cargo de juez de la Corte Suprema es \textbf{vitalicio}. En promedio, los jueces permanecen dieciséis años en funciones.

\section{Phrasal Verbs en Contexto: el caso de ``Make up''}

Un \textbf{phrasal verb} está compuesto por un verbo y una preposición. Su característica principal es que, al combinarse, el significado no siempre es el más obvio que se desprende del verbo y la preposición por separado. 

Es importante identificarlos para no quedarse atrapado tratando de entenderlos por separado. Si se logra identificarlos, se puede partir de la idea de que el significado se lo da el conjunto de las dos palabras, y que muy probablemente no sea el significado más transparente o literal que se desprende de ese verbo y esa preposición.

\subsection*{Distintos significados de ``make up''}

\begin{table}[H]
\centering
\begin{tabularx}{\textwidth}{p{3cm} X X}
\toprule
\textbf{Expresión} & \textbf{Significado} & \textbf{Ejemplo} \\
\midrule
\textbf{make (something) up} & Inventar. & ``Why didn't you call me yesterday?'' --- ``I made it up'' (inventé cualquier cosa). \\
\midrule
\textbf{makeup / to make up} & Maquillaje / maquillarse. & Ponerse rubor, colores, etc. \\
\midrule
\textbf{make up for (something)} & Compensar un perjuicio o desventaja. & Compensar algo que causó un daño. \\
\midrule
\textbf{to make up (pareja o amigos)} & Reconciliarse, recomponer una relación. & Una pareja se pelea, no se habla, y luego ``they make up''. \\
\midrule
\textbf{makeup exam} & Examen recuperatorio. & Recuperar una evaluación en la que salió mal. \\
\bottomrule
\end{tabularx}
\end{table}

El contexto lo es todo. Fíjese que los phrasal verbs son bastante informales; es raro que en un texto de nivel más elevado, o en una sentencia judicial, aparezcan tanto. Son más propios del lenguaje informal. El texto que se está trabajando en clase procede de una página gubernamental, dirigida al público general y no a especialistas en derecho, por lo cual utiliza un lenguaje asequible y comprensible para toda la población, explicando la presencia de estas estructuras informales.

\newpage

%====== PARTE 3: TRIBUNALES FEDERALES ======
\chapter{El Sistema de Tribunales Federales}

\section{La Creación de Tribunales Federales por el Congreso}

La Constitución de los Estados Unidos otorga al Congreso la autoridad para establecer otros tribunales federales, más allá de la Corte Suprema ya constituida.

\begin{important}
El Congreso tiene la potestad de establecer otros tribunales federales para \textbf{atender causas vinculadas con leyes federales}, incluyendo:
\begin{itemize}
    \item Leyes impositivas
    \item Leyes de concursos y quiebras
    \item Causas donde sea parte el Estado federal o los gobiernos de cada uno de los estados
    \item Causas que involucren la Constitución
    \item Y otras materias federales
\end{itemize}
\end{important}

\subsection{El problema de ``state''}

En inglés, \textit{state} es una palabra que para la traducción al español resulta conflictiva. Si se la traduce literalmente como ``estado'', se genera confusión porque ``estado'' también se usa para traducir ``government''. En Argentina no existe este problema porque se distingue entre ``estatal'' (federal) y ``provincial'' (de cada provincia).

Hace algunos años, un grupo de traductores intentó instalar el término ``estadual'' (para los gobiernos de cada estado) frente a ``estatal'' (para el gobierno federal). La propuesta no prosperó porque el término no existe en español y tampoco resultaba claro. El problema no es de la palabra \textit{state} en sí, sino de cómo trasladar esa realidad político-administrativa a nuestro idioma, donde la subdivisión territorial es provincial y no estatal.

\section{Mapa de Tribunales y Agencias Judiciales Especializadas}

El Poder Judicial federal comprende no solo tribunales, sino también organismos que forman parte de su estructura pero no son, estrictamente, tribunales decisorios.

\begin{table}[H]
\centering
\begin{tabularx}{\textwidth}{p{4cm} X}
\toprule
\textbf{Institución (inglés)} & \textbf{Función} \\
\midrule
\textbf{Administrative Office of the US Courts} & Funciones administrativas (no decisorias): presupuesto, estadísticas. \\
\midrule
\textbf{Bankruptcy Courts} & Tribunales con competencia en concursos y quiebras. \\
\midrule
\textbf{Courts of Appeals} & Tribunales de alzada / apelación. \\
\midrule
\textbf{Court of Appeals for the Armed Forces} & Cámara de apelaciones para las fuerzas armadas. \\
\midrule
\textbf{Court of Appeals for the Federal Circuit} & Apelación especializada en patentes y comercio internacional. \\
\midrule
\textbf{Court of Federal Claims} & Demandas contra el Estado (por ejemplo, causas derivadas de privatizaciones). \\
\midrule
\textbf{Court of International Trade} & Comercio internacional. \\
\midrule
\textbf{Federal Court Interpreters} & Servicio de intérpretes judiciales del sistema federal. \\
\midrule
\textbf{Federal Judicial Center} & Investigación y capacitación para jueces. \\
\midrule
\textbf{Judicial Panel on Multidistrict Litigation} & Panel que decide si distintas causas deben consolidarse por tratar una cuestión común (equivalente al ``fuero de atracción''). \\
\midrule
\textbf{US Tax Court} & Tribunal en materia impositiva. \\
\midrule
\textbf{US Court of Appeals for Veteran Claims} & Beneficios y cuestiones de veteranos de guerra. \\
\midrule
\textbf{US Sentencing Commission} & Fija pautas y lineamientos para la imposición de penas, no para su ejecución. \\
\bottomrule
\end{tabularx}
\end{table}

\subsection{La Justicia Militar: una diferencia con el sistema argentino}

Existe una Corte de Apelaciones para las Fuerzas Armadas. En Argentina, en cambio, no hay un tribunal equivalente dentro del organigrama judicial federal. Se cuenta con una Justicia Militar, pero es otro desprendimiento del sistema. Esta diferencia estructural es relevante para cualquier análisis comparado de los sistemas judiciales.

\subsection{El indicador ``US'' en el nombre de un tribunal}

Cuando aparece \textbf{``US''} (United States) adelante del nombre de un tribunal o institución, está marcando que es federal. Esa es la forma de identificar cuándo un tribunal o repartición es federal. Si el nombre no lleva ``US'' e indica en cambio el nombre de un estado (``Court of Appeals for the state of...''), se trata de un tribunal estatal, con su propia organización y denominaciones.

\begin{important}
Es esencial para quienes trabajen en derecho entender quiénes son las partes, cuál es el conflicto, cuál es el tribunal interviniente y cuál es el recorrido de la causa. A veces resulta difícil porque, aunque el sistema es similar al argentino, los nombres son diferentes y la terminología también. Además, cada estado tiene su propio organigrama de cómo está organizada la justicia. Ante la duda, conviene ingresar a la página oficial del Poder Judicial del estado en cuestión para conocer el organigrama exacto de sus instancias.
\end{important}

\section{Denominaciones de los Tribunales: Primera Instancia y Alzada}

En inglés existen distintas maneras de nombrar a los tribunales de primera instancia y de alzada dentro del sistema federal.

\begin{table}[H]
\centering
\begin{tabularx}{\textwidth}{p{3.5cm} p{2.5cm} X}
\toprule
\textbf{Denominación (inglés)} & \textbf{Instancia} & \textbf{Enfoque} \\
\midrule
\textbf{First instance court} & Primera instancia & Nombra directamente la instancia. \\
\midrule
\textbf{District court} & Primera instancia & Nombra la subdivisión territorial (distrito) usada para organizar la justicia federal. \\
\midrule
\textbf{Trial court} & Primera instancia & Nombra la función: es donde se produce la prueba y se lleva adelante el debate. \\
\midrule
\textbf{Court of Appeals / Appellate court} & Alzada & Nombra la función de revisión. \\
\midrule
\textbf{Circuit court} & Alzada & Nombra la subdivisión territorial (circuito). \\
\midrule
\textbf{Upper court / Reviewing court} & Alzada & Nombra la jerarquía (superior) o la función de revisar lo actuado en la instancia inferior. \\
\bottomrule
\end{tabularx}
\end{table}

Todos son distintas formas de referirse a lo mismo. \textit{District} suele traducirse como tribunal de distrito. \textit{Circuit} es una subdivisión territorial a los fines de la administración de justicia; si aparece en el nombre de un tribunal federal, ya se sabe que se trata de alzada. \textit{Trial court} hace referencia a donde se produce la prueba: primera instancia. \textit{Reviewing court} hace referencia a la competencia de apelación, no a la originaria.

\textbf{Recuerde:} ``lower'' y ``upper'' son siempre comparativos relativos. Un tribunal puede ser ``inferior'' respecto de otro sin ser necesariamente de primera instancia.

\newpage

%====== PARTE 4: EL PROCESO JUDICIAL ======
\chapter{El Proceso Judicial: Del Caso al Fallo}

\section{Jurado y Juez: la División de Funciones}

En el sistema anglosajón y estadounidense, los juicios son mayormente por jurado. El jurado resuelve sobre los hechos a partir de la prueba producida; el juez resuelve sobre el derecho.

\begin{table}[H]
\centering
\begin{tabularx}{\textwidth}{p{2.5cm} X X}
\toprule
\textbf{Órgano} & \textbf{Qué resuelve} & \textbf{Cómo se expresa} \\
\midrule
\textbf{Jury (jurado)} & Los hechos, a partir de la prueba producida. & En penal: culpable / no culpable (\textit{guilty / not guilty}). En civil: responsable / no responsable. \\
\midrule
\textbf{Judge (juez)} & El derecho: cómo se aplica la ley a los hechos que fijó el jurado. La cuantía de la pena dentro del máximo y el mínimo legal. & Impone la \textit{sentence} (pena) o resuelve la cuestión de fondo en materia civil (multas, indemnizaciones, medidas restrictivas). \\
\bottomrule
\end{tabularx}
\end{table}

\begin{note}
El jurado no dice ``que se vaya a prisión veinticinco años''; eso lo resuelve el juez, dentro del máximo y el mínimo de la pena. En el sistema argentino, hasta que se popularizó el juicio por jurado, esta función no estaba separada: el juez evaluaba los hechos y después el derecho. En el sistema de juicio por jurado esa función está desdoblada. El jurado no tiene por qué saber de derecho, porque son ciudadanos comunes y corrientes.
\end{note}

Salvo un vicio procesal (por ejemplo, un jurado que recibió dinero o fue objeto de coacción), el juez no puede revisar lo resuelto por el jurado en cuanto a los hechos. Las funciones están bien separadas.

\section{La Etapa Previa al Juicio: Discovery y Hearsay}

\subsection{¿Qué es el Discovery?}

En el sistema estadounidense existe una instancia previa al juicio con jurado, denominada \textbf{Discovery}. En Argentina no existe un equivalente de esa manera.

El \textit{Discovery} funciona como un ofrecimiento de prueba anticipada. Se reciben a testigos que declaran y se anticipa prueba. El juez no está presente físicamente, pero está todo el tiempo: las partes pueden acudir a él ante cualquier anomalía o irregularidad.

Las partes van preparando la estrategia: si un testigo va a decir algo que nos embarra, no lo presentamos cuando estemos delante del jurado, o lo presentamos pero con resguardos. El juez periódicamente convoca a audiencias para monitorear el proceso: ¿cómo van, cuánto les falta? Porque lo que viene después es la etapa con el jurado, salvo que lleguen a un acuerdo antes.

\subsection{Hearsay: Testimonio de Oídas}

\begin{definition}{Hearsay}
Un testigo debe declarar sobre hechos de los que tuvo conocimiento directo. El testimonio basado en lo que un tercero contó (``me dijo la vecina, que le dijo el primo, que le dijo el hermano'') se denomina \textbf{hearsay} y \textbf{no es admisible como prueba}. Si una de las partes intenta presentar un testigo con esas características durante el \textit{Discovery}, la contraparte puede recurrir al juez para impugnarlo.
\end{definition}

También si se presenta un escrito o una foto obtenidos mediante un allanamiento ilegal del domicilio, la contraparte puede acudir al juez para cuestionar esa prueba.

En materia penal, el \textit{Discovery} funciona de otro modo porque es el Estado (a través de la fiscalía) quien impulsa el proceso. En esos casos son frecuentes los acuerdos entre el fiscal y el acusado, por los cuales este último se declara culpable a cambio de una reducción de la pena u otro tipo de beneficio.

\subsection{Sinónimos de ``Juicio'' en Inglés}

El inglés dispone de varios términos, todos traducibles aproximadamente como ``juicio'', ``causa'' o ``proceso'':

\begin{table}[H]
\centering
\begin{tabularx}{\textwidth}{p{2.5cm} X X}
\toprule
\textbf{Término} & \textbf{Uso} & \textbf{Observación} \\
\midrule
\textbf{lawsuit} & El más general para referirse al proceso judicial. & De \textit{sue} (demandar); de uso bastante informal: ``me sumo, hacerme juicio''. \\
\midrule
\textbf{trial} & Proceso en sentido amplio (``to go to trial'': ir a juicio); en sentido más técnico, la etapa probatoria. & Comparable con el doble uso de ``juicio'' en español: a veces todo el proceso, a veces solo la etapa de audiencias. \\
\midrule
\textbf{case} & Causa, caso. & \textit{Case number}: número de causa (expediente). \\
\midrule
\textbf{action / claim} & El acto que da inicio al proceso (``to bring an action''). & No son sinónimos exactos, pero ambos se usan para referirse al juicio en general. \\
\midrule
\textbf{proceeding(s)} & Proceso, en sentido cercano al español. & Abarca las distintas actuaciones dentro de la causa. \\
\bottomrule
\end{tabularx}
\end{table}

Estos términos no son sinónimos exactos entre sí, pero todos se usan para referirse al proceso judicial.

\section{El Camino de Apelación hacia la Corte Suprema}

\subsection{Cómo llega un caso a la Corte Suprema}

No se puede iniciar una demanda directamente ante la Corte Suprema. La mayoría de las causas que llegan a la Corte Suprema lo hacen por apelación.

\begin{table}[H]
\centering
\begin{tabularx}{\textwidth}{p{3.5cm} p{2.5cm} X}
\toprule
\textbf{Vía de acceso} & \textbf{Frecuencia} & \textbf{Cuándo se da} \\
\midrule
\textbf{Appellate jurisdiction} (competencia de apelación) & La mayoría de los casos & Las apelaciones provienen mayormente de tribunales federales, y también de tribunales de los distintos estados cuando el caso trata sobre una cuestión federal (``deals with federal law''). \\
\midrule
\textbf{Original jurisdiction} (competencia originaria) & Ocasionalmente, rara vez & Por ejemplo, en conflictos entre estados, comparables a conflictos entre provincias en el sistema argentino. \\
\bottomrule
\end{tabularx}
\end{table}

\subsection{Requisitos para la Revisión}

Cuando una parte apela, presenta una ``petition'' ante la Corte Suprema, pidiendo que revise la decisión del tribunal inferior. El texto señala: ``Four justices must vote in favor for a case to be granted review''. Se necesitan \textbf{cuatro votos de los nueve jueces} para conceder el recurso y que la Corte acepte revisar la causa.

Los jueces analizan en esta etapa los ``supporting materials'' (memoriales y escritos ya producidos en la instancia anterior), y no nueva prueba (``evidence''), ya que la Corte no recibe elementos probatorios nuevos.

Una vez concedida la revisión, las partes presentan \textbf{``briefs''} (escritos con los fundamentos y argumentos) y argumentos orales. Cada parte dispone de treinta minutos para exponer su caso, durante los cuales los jueces pueden intervenir y repreguntar. Treinta minutos es muy poco para una causa compleja.

\subsection{Las Opiniones Judicales: Mayoría y Disidencia}

Tras los argumentos orales, cada juez elabora su voto (``opinion'').

\begin{table}[H]
\centering
\begin{tabularx}{\textwidth}{p{2.5cm} X X}
\toprule
\textbf{Concepto} & \textbf{Traducción} & \textbf{Explicación} \\
\midrule
\textbf{Opinion} & Voto & Fundamentación individual de cada juez. \\
\midrule
\textbf{Majority opinion} & Voto de la mayoría & Cuando más de la mitad de los jueces vota en el mismo sentido, se convierte en la decisión de la Corte. \\
\midrule
\textbf{Dissenting / minority opinion} & Voto en disidencia / voto minoritario & El voto de quienes no acompañan a la mayoría. \\
\midrule
\textbf{Hands down (a decision)} & Dictar sentencia & Una vez que todos presentaron sus votos y se emitió la resolución, esta ya no puede modificarse. \\
\bottomrule
\end{tabularx}
\end{table}

Es importante comprender que la resolución de la Corte no necesariamente coincide con el primer voto que se lee. Cada juez presenta su voto, y la mayoría de los votos es la que se convierte en la decisión oficial. Mientras los jueces elaboran sus votos debe existir intercambio entre ellos. Como tribunal colegiado, es razonable suponer que en la elaboración de los votos hay intercambios, y que en algún momento uno le pueda decir a otro: ``tenés razón con esto que planteás'', y a lo mejor da vuelta su voto.

\subsection{Volumen de Solicitudes y Tiempos Procesales}

La Corte Suprema recibe entre 7000 y 8000 solicitudes de revisión por año y concede entre 70 y 80. La proporción es muy baja, lo que refleja la selectividad del tribunal.

\begin{note}
La lentitud de los procesos judiciales es un problema estructural. Las causas pueden permanecer en el sistema durante años, acumulando prueba sucesivamente, lo cual genera enormes costos para los litigantes y también para la administración de justicia.
\end{note}

\section{Designación de los Jueces y Comparación con el Sistema Argentino}

\subsection{Mecanismo de Designación}

El proceso de designación de un nuevo juez de la Corte Suprema comprende dos pasos:

\begin{enumerate}
    \item \textbf{Nominación:} El presidente propone un candidato, habitualmente un juez federal con una carrera previa dentro del sistema.
    \item \textbf{Confirmación:} El Senado vota para confirmar la propuesta del presidente.
\end{enumerate}

El cargo es vitalicio. En promedio, los jueces permanecen dieciséis años en funciones.

\begin{important}
Un dato relevante: la Corte puede seguir trabajando y resolviendo causas con menos de nueve jueces; pero en caso de empate no puede revocar la decisión del tribunal inferior, que queda firme.
\end{important}

\subsection{Bench Trial vs. Trial by Jury}

Existe el ``trial by jury'' (juicio por jurado), pero también el ``bench trial'': un juicio en el que no hay jurado, y quien lo resuelve es el juez. No todos los juicios son por jurado. El término \textit{bench} (estrado, banca) proviene de un mobiliario antiguo y muy formal, y se terminó llamando así a ese tipo de juicios.

\subsection{Formalidades del Sistema Anglosajón}

Aunque el derecho anglosajón y el idioma inglés son en algunos aspectos más informales, siguen manteniendo formalidades muy marcadas. En Estados Unidos se usan togas; en el Reino Unido además usan peluca, lo cual no se vive como algo ridículo sino como una formalidad, algo de integración que hay que usar. En España tampoco usan peluca, pero también usan togas.

\subsection{Comparación Final: Estados Unidos vs. Argentina}

\begin{table}[H]
\centering
\begin{tabularx}{\textwidth}{p{3.5cm} X X}
\toprule
\textbf{Aspecto} & \textbf{Estados Unidos} & \textbf{Argentina} \\
\midrule
\textbf{Leyes de fondo y de forma} & Cada estado tiene su propia organización judicial y sus propias leyes. Un mismo hecho puede ser delito en un estado y no en otro. & Códigos de fondo únicos para todo el país (Código Civil y Comercial, Código Penal). Las leyes procesales sí varían según cada provincia. \\
\midrule
\textbf{Pena de muerte} & Existe en algunos estados y en otros no. & No está prevista en el ordenamiento vigente. \\
\midrule
\textbf{Juicio por jurado} & Es la regla en el sistema estadounidense. El jurado resuelve los hechos; el juez, el derecho. & Tradicionalmente el juez resolvía tanto hechos como derecho. El juicio por jurado se encuentra en un proceso de implementación gradual. \\
\midrule
\textbf{Justicia militar} & No existe un tribunal específico dentro del organigrama judicial federal para las fuerzas armadas (más allá de la Corte de Apelaciones para las Fuerzas Armadas). & Existe un fuero militar diferenciado, cuya intervención en delitos de otra naturaleza puede resultar controvertida. \\
\bottomrule
\end{tabularx}
\end{table}

\newpage

%====== CUADRO CORNELL ======
\chapter*{Resumen Cornell}

\section*{Notas y Preguntas Clave}

\begin{table}[H]
\centering
\begin{tabularx}{\textwidth}{
    >{\raggedright\arraybackslash}p{0.25\textwidth}
    >{\raggedright\arraybackslash}X
}
\toprule
\textbf{\small Preguntas / Palabras Clave} &
\textbf{\small Notas / Respuestas} \\
\midrule

\textbf{¿Por qué ``law'' puede significar cosas tan distintas?} &
Porque es un término polisémico: según el contexto puede referirse a (1) una norma en particular, (2) un sistema jurídico en su conjunto, (3) un área o rama del derecho, o (4) la profesión legal (jueces, abogados, fiscales) e incluso, informalmente, a la policía. \\

\midrule

\textbf{¿``Court'' siempre se traduce como ``corte''?} &
No. ``Court'' es el término genérico para tribunal; ``corte'' solo corresponde a la Corte Suprema. Un tribunal de alzada es una cámara, y uno de primera instancia es un juzgado o tribunal. Traducir siempre ``corte'' es un error frecuente. \\

\midrule

\textbf{¿``Sentence'' es lo mismo que ``sentencia''?} &
No. ``Sentence'' es la pena impuesta a quien ya fue declarado culpable; solo se usa cuando hubo condena. ``Sentencia'' en sentido amplio (condenatoria o absolutoria) se traduce como \textit{judgement} o \textit{resolution}. \\

\midrule

\textbf{¿``Jurisdiction'' equivale siempre a ``jurisdicción''?} &
En la práctica, la mayoría de las veces ``jurisdiction'' se traduce como ``competencia'' (territorial, por materia, por grado). ``Jurisdicción'' en sentido estricto (la potestad de todo juez de decir el derecho) se reserva para textos de doctrina. \\

\midrule

\textbf{¿Qué diferencia hay entre comparativo y superlativo?} &
El comparativo compara un elemento con otro. El superlativo lo compara con todo el universo de elementos similares. La Corte Suprema es ``the highest court'' porque no hay ningún tribunal por encima de ella. \\

\midrule

\textbf{¿Qué indican los verbos modales must, may, might, should, shall?} &
Must: obligación. May: permiso o posibilidad. Might: probabilidad remota. Should: recomendación, no obligación. Shall: la obligación por excelencia, propia del lenguaje jurídico y contractual. \\

\midrule

\textbf{¿Cómo está compuesta la Corte Suprema?} &
Nueve miembros: un Chief Justice (presidente) y ocho Associate Justices (jueces asociados). Se necesita un quórum de seis jueces para decidir un caso. \\

\midrule

\textbf{¿Cuánto dura el mandato de un juez de la Corte Suprema?} &
No tiene plazo fijo: es vitalicio, salvo jubilación, renuncia, fallecimiento o remoción por causas establecidas. En promedio, permanecen dieciséis años en el cargo. \\

\midrule

\textbf{¿Qué es un phrasal verb?} &
Combinación de verbo y preposición cuyo significado conjunto no siempre se deduce de sus partes (ej. \textit{make up}: inventar, maquillar, compensar, reconciliarse, recuperar un examen). Identificarlo evita traducir literalmente palabra por palabra. \\

\midrule

\textbf{¿Quién crea nuevos tribunales federales?} &
La Constitución le otorga esa potestad al Congreso, para atender causas vinculadas con leyes federales (impositivas, concursos y quiebras, causas donde sea parte el Estado federal o los estados, casos que involucren la Constitución). \\

\midrule

\textbf{¿Qué indica el prefijo ``US''?} &
Que se trata de un tribunal federal. Si el nombre indica el nombre de un estado, se trata de un tribunal estatal. \\

\midrule

\textbf{¿Qué diferencia hay entre trial court, district court y court of appeals?} &
\textit{Trial court} y \textit{district court} son formas distintas de nombrar tribunales de primera instancia. \textit{Court of Appeals}, \textit{circuit court} y \textit{reviewing court} nombran tribunales de alzada. \\

\midrule

\textbf{¿Qué resuelve el jurado y qué el juez?} &
El jurado resuelve los hechos (culpable/no culpable; responsable/no responsable). El juez resuelve el derecho y determina la pena dentro del máximo y el mínimo legal. \\

\midrule

\textbf{¿Qué es el Discovery y el hearsay?} &
\textit{Discovery} es la etapa previa al juicio donde se anticipa prueba. \textit{Hearsay} es el testimonio de oídas, no admisible como prueba. \\

\midrule

\textbf{¿Cómo llega un caso a la Corte Suprema?} &
Mayormente por apelación (necesita cuatro votos de nueve jueces). Rara vez por competencia originaria, como en conflictos entre estados. \\

\midrule

\textbf{¿Qué es la majority opinion y la dissenting opinion?} &
La \textit{majority opinion} es el voto de más de la mitad de los jueces, que se convierte en la decisión de la Corte. La \textit{dissenting opinion} es el voto de la minoría. \\

\midrule

\textbf{¿Cómo se designa a un juez de la Corte Suprema?} &
El presidente propone un candidato; el Senado vota para confirmar. El cargo es vitalicio. Con empate, no puede revocarse la decisión inferior, que queda firme. \\

\midrule \midrule

\multicolumn{2}{p{\textwidth}}{\textbf{Síntesis Final:} La clase reconstruye la organización del Poder Judicial federal estadounidense (Corte Suprema, tribunales inferiores, agencias especializadas) a partir de un texto oficial en inglés. Se entrena lectocomprensión mediante la identificación de polisemia, falsos cognados, verbos modales y phrasal verbs, con comparación constante con el sistema judicial argentino. Se destacan diferencias estructurales clave: división de funciones entre jurado y juez, sistemas judiciales estatales independientes, y procedimientos como el \textit{Discovery}.} \\

\bottomrule
\end{tabularx}
\end{table}

\end{document}
